\documentclass[11pt,a4paper]{article}
\usepackage[margin=27mm]{geometry}
\usepackage{fontspec,unicode-math,amsmath,amsthm,booktabs,longtable,hyperref}
\setmainfont{STIXGeneral}\setmathfont{STIX Two Math}\hypersetup{hidelinks}
\newtheorem{theorem}{Theorem}\newtheorem{lemma}[theorem]{Lemma}
\newtheorem{proposition}[theorem]{Proposition}
\newcommand{\PSL}{\operatorname{PSL}}\newcommand{\Aut}{\operatorname{Aut}}
\newcommand{\Cl}{\operatorname{Cl}}\newcommand{\Cen}{\operatorname{C}}
\title{The universal locally projective polytope and its corrected quotients}
\author{Application to the Hartley--Leemans 2003 preprint}\date{}
\begin{document}\maketitle
We recover the complete corrected principal contract of the fifteen-page arXiv edition \cite{source}, with the later two-page erratum \cite{erratum} kept as a distinct original. The common paper \cite{common} supplies all marked-group, geometric and exhaustive-census proofs, its entire five-page table companion and all six required programs. This application states their consequences in the original notation and preserves the exact presentation, covering scope and source qualifications. It does not assume the other application to the erratum.

\section{Literal presentation and finite structure}
Let $W=\langle s_0,s_1,s_2,s_3\rangle$ have the relations
\[
 s_i^2=1,\quad(s_0s_1)^5=(s_1s_2)^3=(s_2s_3)^5=1,
 \quad[s_i,s_j]=1\ (|i-j|>1),\quad(s_1s_2s_3)^5=1,
\]
and set $\omega=(s_0s_1s_2)^5$, $\nu=\omega s_3$.
\begin{theorem}\label{group}
The literal presented group is $J_1\times\PSL_2(19)$ of order $600415200$. The normal closures of $\nu^3$ and $\omega$ are respectively the projective and Janko factors, giving the exact marked quotient presentations of $J_1$ and $\PSL_2(19)$. The element $\nu$ has order six. The source subgroup
\[
 L=\langle s_2s_1,\ (s_0s_1)^2s_2s_3s_0,\ s_0s_1s_0s_3s_2s_3s_1\rangle
\]
has index $20$, order $30020760$, and its specified source generators $x,y$ satisfy the exact presentation
\[
 L=\langle x,y\mid y^3=xyxy^{-1}xyx^{-1}yx=1\rangle.
\]
The group has the faithful direct-product action on $266+20=286$ points, and its displayed twenty-coset action has the Janko factor as kernel.
\end{theorem}
\begin{proof}
Apply Theorem 6 of the common paper. Its proof starts in the actual presentation: complete enumeration gives the index $20$, and tracked rewriting identifies the source words $x=u_2^{-1}u_3u_1$, $y=u_1u_3^{-1}$ with the two resulting generators. The cyclic subgroup on $x$ has index $175560$ and a complete one-column relation matrix of gcd $171$. Lemma 4 of the common paper bounds $|W|$ by $20\cdot175560\cdot171$. Only then is the finite product image of that same order used for equality. Thus the displayed finite action is not mistaken for an upper-bound proof. The checked coordinate orders and simplicity establish the exact normal closures and order of $\nu$; the literal twenty source coset words and permutations are included in the complete program data. The two-generator presentation gives $|L|=175560\cdot171=30020760$.
\end{proof}

\section{Universality and the distinguished simple quotients}
\begin{theorem}\label{geometry}
The above marking is a string C-group. Its polytope $\mathcal P$ is universal for dodecahedral facets and hemi-icosahedral vertex figures. The Janko quotient has $1463$ dodecahedral facets and $2926$ vertices. The projective quotient is the self-dual $57$-cell with hemi-dodecahedral facets and hemi-icosahedral vertex figures. Neither distinguished quotient has a proper polytope quotient in the semisparse sense. Every section-regular locally projective polytope of type $\{5,3,5\}$, or its dual, occurs in the resulting quotient classification. The complete classification also includes the non-section-regular mixed-facet quotients.
\end{theorem}
\begin{proof}
Theorem 7 of the common paper checks all $256$ subset intersections for each of the three actual markings, the exact adjacent orders and the nonadjacent commutations. It identifies the endpoint groups by their marked classical presentations and orders, rather than by an abstract isomorphism alone. The endpoint orders are $(120,60)$ for $W$ and $J_1$, and $(60,60)$ for the projective quotient. Their indices give the stated face counts. Reversal of the actual projective markers extends to a bijective homomorphism and proves self-duality.

The same common theorem proves universality using the flag-adjacency action of the literal presentation; it does not assume that the target is regular. The local projective alternatives and duality then give the stated section-regular scope. Hartley's retained general endpoint theorem reduces the two proper-quotient exclusions to the actual forbidden endpoint products. Lemma 5 of the common paper and the complete prime-order subgroup tests exclude every nontrivial semisparse subgroup. Its proof also checks all five literal Janko witness words of orders $3,5,19,11,7$. The projective test uses prime-order subgroup classes, not the false assumption that a Sylow $3$-subgroup has order three.
\end{proof}

\section{The complete corrected quotient classification}
\begin{theorem}\label{census}
The full quotient classification for $\mathcal P$ has $137$ isomorphism classes, with $26$ maximal semisparse subgroup classes and $349$ maximal-containment edges. It contains $65$ quotients with the original local types, one with only hemi-dodecahedral facets, and $71$ with mixed facets. Exactly $66$ are section regular. All facet counts, all actual subgroup representatives and all quotient automorphism groups are specified in the complete common census tables.
\end{theorem}
\begin{proof}
Theorem 8 of the common paper uses the already established marked group and geometry, the exact endpoint criterion, and its Lemmas 1 and 3 and Proposition 2. It tests all $1262$ subgroup conjugacy classes and reconciles all $145$ literal historical rows. The $137$ survivors, all containment edges, facet counts and normalizer quotients are the output of the complete required calculation. The eight deleted rows are $13,70,78,90,91,95,111,127$, each with an explicit forbidden-conjugacy witness. In particular no erroneous historical total is a premise, and no sampling claim is substituted for exhaustiveness.
\end{proof}
The entire five-page table companion to the common paper is part of this result, including both full tables and every surviving row. For example row $143$ has $3420$ hemi-facets and automorphism group $A_5\times\PSL_2(19)$; the full rows retain the ordinary facets as well. The thirty source maximal-subgroup word lists are evaluated in the common code; $26$ survive. Downward closure follows from the complete conjugacy test, and the retained general quotient correspondence turns subgroup inclusion into covering.

\section{Editions, actions and attribution}
The accessible source is arXiv:math/0307075v1, fifteen pages. We do not identify it with the unseen twelve-page journal edition. The corrected principal totals are $137,26,65,66,71$ in place of the original $145,30,69,70,75$. The erratum remains a distinct two-page original, although both applications reuse the same complete common proof and table output.

The arXiv Table 1 order-$70$ word labelled $58$ maps to full-table row $63$; full-table row $58$ instead has order $76$. Row $21$ prints $2^2\times(19:3)$ although its actual order is $456$ and its group is $2^3\times(19:3)$. Row $27$ prints $2\times19:3$ although its order is $342$ and its group is $2\times(19:9)$. Row $41$ has elementary abelian normal subgroup of order eight and exponent two, so its printed $8:(7:3)$ is not a cyclic-eight action under the source's stated notation. Rows $38$--$40$ preserve different actions: their actual centers have orders $1,1,3$. Abbreviated structure strings and orders are not used alone to prove subgroup isomorphism. Exact permutation representatives preserve the actions, and every normalizer quotient is checked by the complete common reconciliation.

The universal-group and geometric results, source words and historical tables are Hartley--Leemans results. The marking search and exhaustive corrected census are previously reviewed work reused in full. The present modular edition introduces no new finite calculation. Its common proof includes the presentation upper bound and all local adapters instead of relying on the missing historical \texttt{g.txt}. The old forty-day runtime, software history and acknowledgements remain historical reports, not mathematical assertions to reproduce. Classical string-C-group and semisparse correspondences, Hartley's general endpoint criterion, local marked spherical and antipodal groups, and named GAP, ACE, TomLib and NetworkX algorithms remain explicit retained inputs. No complete proof cost for those external inputs is asserted, and there is no new Lean declaration, new question solution or priority claim.
\begin{thebibliography}{9}
\bibitem{common} \emph{Marked groups and complete locally projective quotient censuses}, common paper, complete five-page census companion and six whole required programs in the present edition.
\bibitem{source} M. I. Hartley and D. Leemans, \emph{Quotients of a Universal Locally Projective Polytope of type $\{5,3,5\}$}, arXiv:math/0307075v1, 5 July 2003, fifteen pages. Journal reference: Math. Z. \textbf{247} (2004), 663--674; the journal edition has not been collated here.
\bibitem{erratum} M. I. Hartley and D. Leemans, erratum, Math. Z. \textbf{253} (2006), 433--434, doi:10.1007/s00209-005-0872-z.
\end{thebibliography}
\end{document}
